\section*{第 \ref{ch:g4:plants-without-seeds} 章——无种子植物}

\begin{solution}{exo:g4:plants-without-seeds:1}
\omterm{prop:g4:plants-without-seeds:damp}{苔藓}和\omterm{prop:g4:plants-without-seeds:damp}{蕨类}。它们靠\omterm{def:g4:plants-without-seeds:spore}{孢子}繁殖，\omterm{def:g4:plants-without-seeds:spore}{孢子}是
数量极多、微小的颗粒，被大量释放。
\end{solution}

\begin{solution}{exo:g4:plants-without-seeds:2}
\omterm{def:g4:plants-without-seeds:spore}{孢子}微小而裸露——没有食物储备，也没有
现成的微型植株；\omterm{def:g2:from-seed-to-plant:seed}{种子}两者兼备，外面包着
保护性的种皮。而且\omterm{def:g4:plants-without-seeds:spore}{孢子}以百万计产生，\omterm{def:g2:from-seed-to-plant:seed}{种子}却只产生十二个左右。
\end{solution}

\begin{solution}{exo:g4:plants-without-seeds:3}
在蕨叶背面，排列整齐的褐色小点，每个
都是一个装有数千个\omterm{def:g4:plants-without-seeds:spore}{孢子}的小囊，成熟时裂开，
让\omterm{def:g4:plants-without-seeds:spore}{孢子}像尘埃一样飘散而去。
\end{solution}

\begin{solution}{exo:g4:plants-without-seeds:4}
\omterm{prop:g4:plants-without-seeds:damp}{苔藓}靠\omterm{def:g4:plants-without-seeds:spore}{孢子}繁殖，而\omterm{def:g4:plants-without-seeds:spore}{孢子}及其娇嫩的
幼体阶段都需要潮湿——在干燥的空气中，它们会失水死亡。墙的北侧
背阴处始终潮湿，南侧向阳的墙面则被烈日烤干。
\end{solution}

\begin{solution}{exo:g4:plants-without-seeds:5}
直接由亲代植株的一部分长出新植株
——如\omterm{ex:g4:plants-without-seeds:runner}{匍匐茎}、\omterm{ex:g4:plants-without-seeds:bulbtuber}{鳞茎}、\omterm{ex:g4:plants-without-seeds:bulbtuber}{块茎}和\omterm{met:g4:plants-without-seeds:cutting}{扦插}。只有一个亲本。
\end{solution}

\begin{solution}{exo:g4:plants-without-seeds:6}
植株抽出一根细长的茎——即\omterm{ex:g4:plants-without-seeds:runner}{匍匐茎}——
横卧在土壤上；茎与土壤接触处生根、长出叶，
形成一株完整的小植株；随后\omterm{ex:g4:plants-without-seeds:runner}{匍匐茎}枯萎，
子株便独立出来，与亲本完全相同。
\end{solution}

\begin{solution}{exo:g4:plants-without-seeds:7}
种下的马铃薯（即\omterm{ex:g4:plants-without-seeds:bulbtuber}{块茎}）能从其
“芽眼”可靠地长出植株，每个芽眼都能长成一整株，
与所种的优良品种完全相同——比\omterm{def:g2:from-seed-to-plant:seed}{种子}更快、更有把握。
\end{solution}

\begin{solution}{exo:g4:plants-without-seeds:8}
答案不唯一。白色的根从浸没的茎上长出，
通常在一至两周后形成；此时\omterm{met:g4:plants-without-seeds:cutting}{扦插}枝即可
栽种——从一个亲本得到一株完整的新薄荷。
\end{solution}

\begin{solution}{exo:g4:plants-without-seeds:9}
\omterm{def:g4:plants-without-seeds:spore}{孢子}裸露而成本低廉：没有食物储备，也没有
小植株，几乎每一个都会丢失，所以必须产生数以百万计。\omterm{def:g2:from-seed-to-plant:seed}{种子}
则是昂贵的生存装备，十二个装备精良的\omterm{def:g2:from-seed-to-plant:seed}{种子}便已足够。它们与
鳕鱼（数以百万计，无人照料）和
大象（数量少，照料周到）形成呼应。
\end{solution}

\begin{solution}{exo:g4:plants-without-seeds:10}
草莓是借\omterm{ex:g4:plants-without-seeds:runner}{匍匐茎}繁殖的——即\omterm{def:g4:plants-without-seeds:vegetative}{营养繁殖}，
只有一个亲本，得到的是完全相同的复制品。豆子是由花产生的\omterm{def:g2:from-seed-to-plant:seed}{种子}长成的，
其中混合了两个亲本的贡献（\omterm{def:g3:flowers-fruits-seeds:parts}{花粉}来自一朵花，
\omterm{def:g4:animal-reproduction:sexual}{卵细胞}来自另一朵花），因此每株豆苗
都略有差异。
\end{solution}

\begin{solution}{exo:g4:plants-without-seeds:11}
\omterm{def:g4:plants-without-seeds:vegetative}{营养繁殖}——这根枝条本身就是天然的\omterm{met:g4:plants-without-seeds:cutting}{扦插}材料。
洪水代替园丁完成了剪枝的工作：它切下枝条，
将其冲走，并栽进湿润的泥土中，那正是柳树生根的理想场所。
河流两岸的柳树，正是这样在同一瞬间
完成折断与栽种而长成的。
\end{solution}
